%! Unit Launch: Study Design — Unit 1, Lesson 1.0 — Warm-Up (Answer Key)
\documentclass[10pt]{article}
\usepackage{saar-article}
\usepackage{saar-key}   % pulls in -boxes; do NOT also load -boxes
\newcommand{\TallMath}[1]{$\displaystyle #1\rule[-1.4em]{0pt}{3.2em}$}

\begin{document}

\pageheader{Unit 1, Lesson 1.0}{Warm-Up --- Answer Key}
% NAMESTRIP: no \namedateperiod here — mirrors the blank. Teacher-only prose goes
% in the lesson plan as \begin{teachernote}[Warm-Up], never in a key.

\vspace{0.15in}

\begin{notesbox}{Spiral Review --- Percents, Proportions, and the Mean}
\small
Everything below returns this unit. Show your work.

\begin{enumerate}[leftmargin=1.5em, itemsep=6pt, topsep=2pt]

  \item Riverbend High School has $240$ seniors. Of those seniors, $84$ ride the
        bus to school. What \textbf{percent} of the seniors ride the bus?

        \begin{work}
          \dfrac{84}{240} &= 0.35 \\
                          &= 35\%
        \end{work}

  \item A survey of $50$ Riverbend students found that $18$ of them walk to
        school. Riverbend has $900$ students in all. Based on this survey,
        \textbf{predict} how many of the $900$ students walk to school.

        \begin{work}
          \dfrac{18}{50} &= 0.36 \\
          0.36 \times 900 &= 324 \text{ students}
        \end{work}

  \item Five students recorded how many minutes they spent on homework last
        night: $12$, $15$, $9$, $20$, $14$.

        \textbf{(a)} Find the \textbf{mean} number of minutes.

        \begin{work}
          \bar{x} &= \dfrac{12+15+9+20+14}{5} \\
                  &= \dfrac{70}{5} \\
                  &= 14 \text{ minutes}
        \end{work}

        \textbf{(b)} Write one sentence that says what your answer in part (a)
        describes. Name the group and the units.
        \ansline{The five students surveyed spent a mean of 14 minutes on homework last night.}

\end{enumerate}
\end{notesbox}

\end{document}
